%%%%%%%%%%%%%%%%%%%%%%%%%%%%%%%%%%%%%%%%%%%%%%%%%%%
% 姜睿 | 端侧视觉算法 / 深度模型部署 / 影像质量工程
%%%%%%%%%%%%%%%%%%%%%%%%%%%%%%%%%%%%%%%%%%%%%%%%%%%
\name{姜睿}
\contactInfo{17726952709}{Jru0614@163.com}{甘肃兰州}

\logosection{\faGraduationCap}{教育经历}
\datedline{\textbf{清华大学｜电子信息｜硕士}}{2025.09 -- 2027.06}
\datedline{\textbf{哈尔滨工业大学｜精密仪器｜本科}}{2020.09 -- 2024.06}
\begin{itemize}
  \item 五四表彰优秀学生、两次人民奖学金；大创国家级一等奖。
  \item 第八届智能车竞赛智能视觉组北部赛区三等奖；第九届光电设计竞赛东北赛区二等奖。
\end{itemize}

\logosection{\faSuitcase}{实习经历与项目}
\datedline{\textbf{vivo｜影像效果工程师}}{2026.06 -- 2026.08}
\datedline{\textbf{LiveAgent｜智能模板生成与端侧 SDK}}{}
\begin{itemize}
  \item \textbf{背景与任务}：面向 Live 相册的空间拼贴玩法，参与搭建从拼贴样例中解析模板并复用的统一框架，降低新玩法对人工编排和重复开发的依赖。
  \item \textbf{模板解析与复用}：以一张空间拼接样例作为输入，由 VLM\textbf{“大脑”}理解画面中的布局、素材关系和组合结构，并调用图像分割、目标追踪等模型提取素材区域及其空间约束，解析生成可复用的拼贴模板。
  \item \textbf{素材理解与布局}：用户仅需选择待组合素材，大脑模型根据素材内容理解其主体和语义特征，匹配模板中的位置与关系并完成自动布置，支持不同素材在同一模板下的灵活替换和组合。
  \item \textbf{玩法扩展与端侧落地}：在拼贴结果上集成滤镜变化、文本添加及基于 DiT 的图像变换玩法；将原型架构翻译为 Kotlin 并封装为手机端 SDK，完成模块接口、数据传递和端侧调用联调，当前已实现核心功能闭环，性能指标待进一步测试。
\end{itemize}

\datedline{\textbf{算能科技｜芯片规划工程师}}{2026.01 -- 2026.05}
\datedline{\textbf{深度模型部署与端侧推理优化}}{}
\begin{itemize}
  \item \textbf{背景与任务}：面向 CV186 Mars3（BM1684X、UMA 架构）端侧部署 300M 参数级单目深度估计模型，针对原生模型在高分辨率实时推理下的显存、带宽和功耗压力，负责模型转换、算子适配和性能优化。
  \item \textbf{计算图优化}：从 PyTorch 导出 ONNX 并定位运行时不支持或效率较低的算子，裁剪 Mask/法线等冗余分支，固化位置编码与 Token 长度；将 `scaled\_dot\_product\_attention` 等价改写为显式 $QK^T$ 计算，提升 NPU 适配性。
  \item \textbf{端侧算子与内存优化}：针对 GroupNorm/LayerNorm、Resample 等模块进行结构改造和精度对齐，在 C++ 前处理节点完成异常值隔离，并手动融合部分拼接操作，减少 UMA 架构下的中间张量和内存搬运。
  \item \textbf{量化与结果}：按节点敏感度设计 FP16/INT8 混合量化，模型参数量由 300M 降至 85M，显存由 1.9GB 降至 0.9GB，带宽由 30GB/s 降至 5GB/s；在 17W 功耗约束下实现 1080P@30fps 稳定推理。
\end{itemize}

\datedline{\textbf{图像质量评估模型与自动化 PQ 调参助手}}{}
\begin{itemize}
  \item \textbf{背景与任务}：面向影像 PQ 调试中“图像差异判断”和“场景因素分析”混杂的问题，参与构建支持 \textbf{JPEG 对比图与 RAW 图} 两类输入的图像质量评估及参数辅助分析流程。
  \item \textbf{VLM 语义解析}：不让 VLM 直接承担像素级画质对比，而是发挥其更擅长的视觉语义理解能力，解析天黑、逆光等环境信息，以及灯光、人像等画面内容组合，并将结果组织为结构化场景标签。
  \item \textbf{知识检索与二次分析}：将场景标签、输入类型和 PQ 问题作为检索条件，从 RAG 知识库及对应 JSON 参数配置中召回 Top-$k$ 候选结果，再结合当前样本进行二次分析，给出更匹配场景的参数调整建议。
  \item \textbf{流程工程化}：参与样本管理、输入解析、模型调用、结果记录和回归验证等环节的串联，形成“样本输入—语义解析—参数检索—建议输出”的可复用调试流程；完成核心功能验证，具体效率指标待后续测试补充。
\end{itemize}

\logosection{\faCogs}{专业技能}
\begin{itemize}
  \item \textbf{视觉与成像基础}：具备相机成像、光学成像链路和图像处理基础；了解 3A、立体视觉、图像防抖及常见影像质量问题。
  \item \textbf{深度学习与 VLM}：Python、PyTorch、Hugging Face Transformers；具备视觉模型推理接入、VLM 语义解析、RAG/JSON 参数检索流程开发经验，了解 DiT 工作流。
  \item \textbf{模型部署与优化}：具备 PyTorch/ONNX 模型导出、计算图改写、算子适配、FP16/INT8 量化、精度验证和端侧性能分析经验；使用过算能 TPU/运行时、QNN 与 MLIR 编译部署工具链。
  \item \textbf{工程开发}：C++、Kotlin、Docker、Git；具备移动端 SDK 接口封装、端侧功能联调、批量测试和日志/结果记录经验；了解基于 MLIR 的计算图转换、算子 lowering 与端侧编译流程。
  \item \textbf{PQ 与工具链}：了解 JPEG/RAW 图像处理链路，具备 PQ 参数调试、样本管理、回归验证及自动化调参流程搭建经验。
  \item \textbf{其他信息}：英语 CET-6；可实习半年以上。
\end{itemize}
%%%%%%%%%%%%%%%%%%%%%%%%%%%%%%%%%%%%%%%%%%%%%%%%%%%
